\begin{notation}
  Seja $V=\mathbb{R}^{n}$.\\
  Neste caso denotamos
  \begin{align*}
    e^{i}=dx^{i}.
  \end{align*}
\end{notation}
\begin{definition}[Forma diferencial]
  Uma \textbf{forma diferencial} de grau $k$ em um aberto $U\subseteq\mathbb{R}^{n}$ é uma função
  \begin{align*}
    w:U\to\Lambda^{k}(\mathbb{R}^{n}).
  \end{align*}
  diferenciável no seguinte sentido, escrevendo em coordenadas locais, pegando $I$ multi-índice crescente 
  \begin{align*}
    w(x)=\sum_{I}w_{i_1,\cdots,i_k}(x)dx^{i_1}\wedge\cdots\wedge dx^{i_k}=\sum_{I}w_{I}dx^{I}.
  \end{align*}
  Onde as funções componentes
  \begin{align*}
    w_{I}:U\to\mathbb{R}.
  \end{align*}
  são todas diferenciaveis.\\
  A forma diferenciável é de classe $C^{m}$ se todas as funções componentes são $C^{m}$. 
\end{definition}
\begin{note}
  \begin{itemize}
    \item Note que
      \begin{align*}
        w+\eta=&\sum_{I}(w_I+\eta_I)dx^{I}.\\
        \lambda w=&\sum_{I}\lambda w_Idx^{I}.\\
        w\wedge\eta=\sum_{IJ}w_I\eta_{J}dx^{I}\wedge dx^{J}.
      \end{align*}
    \item Se $\Lambda^{k}(\mathbb{R}^{n})$ é o espaço vetorial das $k$-formas diferenciáveis, temos que
    \begin{align*}
      \Lambda(\mathbb{R}^{n})=\oplus_{k=0}^{n}\Lambda^{k}(\mathbb{R}^{n})
    \end{align*}
    é a álgebra exterior.
  \end{itemize}
\end{note}
